

\chapter{Service Implementation}

\section{Public Key Infrastructure (PKI)}

\subsection{Architecture Overview}

\textbf{Hierarchy:} Two-tier PKI with Root CA (self-signed, offline) and Intermediate CA (signs all certificates). \textbf{Algorithms:} RSA 4096-bit keys, SHA-256 signatures. \textbf{Validity:} Root CA (10 years), Intermediate CA (5 years), server/client certificates (1 year).

\subsection{Implementation}

\textbf{Directory Structure:} \texttt{PKI/root-ca/} and \texttt{PKI/intermediate-ca/} with \texttt{private/} (600 permissions), \texttt{certs/}, \texttt{crl/}, \texttt{checksums/}, \texttt{scripts/}.

\textbf{Automated Scripts:}
\begin{itemize}
    \item \texttt{init-ca.sh}: Initialize CA databases (index.txt, serial)
    \item \texttt{generate-root-ca.sh}: Create Root CA (RSA 4096, SHA-256, 10-year validity)
    \item \texttt{generate-intermediate-ca.sh}: Create Intermediate CA signed by Root CA
    \item \texttt{issue-server-cert.sh <hostname>}: Issue server certificates with \texttt{extendedKeyUsage = serverAuth}
    \item \texttt{issue-client-cert.sh <name>}: Issue client certificates with \texttt{extendedKeyUsage = clientAuth}
    \item \texttt{revoke-cert.sh <cert>}: Revoke certificates and update CRL
    \item \texttt{verify-integrity.sh}: Verify SHA-256 checksums of critical files
\end{itemize}

\textbf{Security Features:}
\begin{itemize}
    \item Dedicated user \texttt{ca} with restricted permissions (private keys 600, directories 700)
    \item SHA-256 checksums for all critical files (keys, configs, database)
    \item CRL generation with 30-day validity (PEM and DER formats)
    \item Certificate chain files included for easy deployment
\end{itemize}

\subsection{Issued Certificates}

\begin{table}[h]
\centering
\begin{tabular}{|l|l|l|}
\hline
\textbf{Certificate} & \textbf{Serial} & \textbf{Extended Key Usage} \\
\hline
web01.org.local & 1001 & serverAuth (TLS server) \\
web02.org.local & 1004 & serverAuth (TLS server) \\
db01.org.local & 1002 & serverAuth (TLS server) \\
nginx01.org.local & 1005 & serverAuth (TLS server) \\
web02-client & 1006 & clientAuth (mTLS client) \\
\hline
\end{tabular}
\caption{Issued certificates for infrastructure services}
\end{table}

\textbf{Certificate Properties:} RSA 4096-bit, SHA-256 signature, 1-year validity, Subject CN and SAN with hostname, Key Usage (Digital Signature, Key Encipherment), certificate chain files included.

\subsection{Revocation System}

\textbf{CRL Management:} \texttt{revoke-cert.sh} updates database, \texttt{generate-crl.sh} creates CRL (30-day validity, PEM/DER formats), \texttt{verify-revocation.sh} checks revocation status.

\textbf{Test Verification:} Certificate \texttt{test-revoke.org.local} successfully revoked, \texttt{openssl verify -crl\_check} returned \texttt{error 23: certificate revoked}, confirming functional revocation system (Assignment 2 PKI lifecycle requirement).
